\paragraph{Steering makes tokens compete.} Steering enters both factors of \eqref{eq:gain}: it moves the logits, and
with them $\Sigma_p$, and it moves the hidden states, and with them $J_{\tilde h}$. Which change makes the loop
chaotic? Three measurements on Qwen2.5-0.5B, steered along the philosophy direction, give one answer
(Figure~\ref{fig:mechanism}, Table~\ref{tab:mechanism}). First, each step's growth splits exactly into a network
factor $g_{\mathrm{net}}=(|\delta\tilde h|/|\tilde h|)/(|\delta x|/|x|)$ and a readout factor
$g_{\mathrm{read}}=(|\delta F|/|F|)/(|\delta\tilde h|/|\tilde h|)$. Averaged along an orbit, their logarithms add up
to the largest Lyapunov exponent, up to a boundary term. Over a $101\times53$ map in steering strength and temperature
(block 12; 5353 cells, 331 chaotic), the chaotic cells' readout factor exceeds that of nearly unsteered cells
($|\alpha|\le0.01$) at the same temperature by $2.58$ nats per step (95\,\% interval $[1.79,3.58]$ over 18 temperature
rows). Their network factor differs by only $-0.07$ $[-0.21,0.11]$. Without steering the readout contracts strongly,
by $0.9$ to $5.5$ nats per step; in the chaotic cells it expands. Compared with periodic cells at the same temperature,
chaotic cells have a higher next-token entropy ($+1.62$ nats, $[1.26,1.91]$), a smaller gap between the two largest
logits ($-1.37$ in units of $T$, $[-1.98,-0.87]$) and a larger readout factor ($+1.36$, $[0.92,1.92]$), each with
$p\le10^{-4}$ in a permutation test within temperature rows.

Second, a readout-only intervention gives chaos. We add to the logits a bias $\beta\,b$ along the direct logit effect
of the steering vector ($b=Wv$, centred and scaled to unit standard deviation over the vocabulary), with no residual
steering. This leaves the network untouched and changes only which tokens compete. It makes 3.5\,\% of the same
temperature range chaotic (241 of 6837 cells), against 6.2\,\% for residual steering, mostly at
$\log_{10}T>0.1$. Creating token competition therefore suffices for chaos, given the unmodified network.

Third, a gain split compares different orbits, so steering could act simply by moving the orbit to other states. A
crossover at fixed states rules this out. Along recorded steered orbits we evaluated the steered and the unsteered
model at the same states and carried four tangent dynamics. Each pairs the steered or unsteered readout $R(p)$ with the
steered or unsteered network Jacobian: arms 11, 10, 01 and 00. Arm 11 reproduces the twin exponent within $0.012$ at 23
of 24 cells. Over 14 held-out chaotic cells, steering only the readout raises the growth by $0.81$ nats per step
($[0.29,1.35]$). Steering only the network response lowers it by $0.49$ ($[-0.63,-0.35]$). The interaction is $+0.26$
($[0.10,0.45]$). Below $\log_{10}T=0$ the steered readout turns a contracting unsteered derivative ($-1.03$ nats per
step) into an expanding one ($+0.72$), more strongly than the full steered map ($+0.34$). Above it, the unsteered
derivative already expands at the steered states ($+0.41$), and steering acts by bringing the orbit there. In
Qwen2.5-7B, whose embeddings are untied, the same crossover (18 cells of its steered map at block 14; arm 11 equals the
twin exponent within $10^{-4}$ at every cell) reproduces the damping. Over 10 held-out chaotic cells, steering only the
network response lowers the growth by $1.14$ nats per step ($[-1.61,-0.64]$). Eight of these cells lie at
$\log_{10}T\ge0$, where the unsteered derivative already expands at the steered states ($+1.33$), and steering only the
readout adds $+0.23$ ($[-0.11,0.63]$).

Steering thus makes the loop chaotic by making next tokens compete, which turns the softmax readout from a
contraction into an expansion. At low temperature it creates that competition at fixed states. At higher temperature it
mainly carries the orbit into states that are already competitive. In both models its change of the network's response
damps the growth. The readout rule holds below $\log_{10}T\approx0.1$; at higher temperature the network factor itself
can expand. Steered at block 4, the chaotic
cells at $\log_{10}T=0.15$ have a network factor of $+0.72$ and a contracting readout ($-0.35$). Along the hostile line
at $\log_{10}T=0.12$--$0.13$ the network factor of the chaotic cells is $+0.60$ against $+0.24$ to $+0.38$ at the
neighbouring periodic cells. The mechanism is not specific to
concept-labelled directions. This is consistent with the controls of Section~\ref{sec:discussion}, in which
norm-matched random directions also produce chaos in the 0.5B and 8B models; at 70B the directions differ
(Section~\ref{sec:size}). What a direction does select is which tokens compete, and therefore what the attractors
write.
\begin{figure}[tbp]
\centering
\includegraphics[width=\linewidth]{fig-mechanism.png}
\caption{Why steering makes the loop chaotic (Qwen2.5-0.5B-Instruct, philosophy direction, float64).
(a) Readout factor (solid) and network factor (dotted) along steering lines at block 12, at five temperatures; dots mark
chaotic cells. (b) Row-matched differences between chaotic and periodic cells (black) and between chaotic and nearly
unsteered cells (red) on the $101\times53$ map, with 95\,\% bootstrap intervals over temperature rows and
permutation $p$-values. (c) Orbit-preserving crossover: finite-time growth along each recorded steered orbit when the
readout and the network Jacobian are taken from the unsteered (0) or steered (1) model; red, chaotic cells; blue,
periodic cells; dark red, the mean over held-out chaotic cells; navy, the same mean in Qwen2.5-7B. (d) Share of chaotic cells per temperature under
residual steering at block 12 and under a pure logit bias along $Wv$. (e) The two factors of the chaotic cells by
steering block, at $\log_{10}T\le0$. (f) Kaplan--Yorke dimension against the readout's stable rank at the chaotic cells of the dimension
measurement (Section~\ref{sec:theory}); dotted, equality.}
\label{fig:mechanism}
\end{figure}
\begin{table}[!htb]
\centering\small
\begin{tabular}{lll}
\toprule
measurement & contrast & estimate (95\,\% interval) \\
\midrule
gain split, $101\times53$ map & readout factor, chaotic $-$ $|\alpha|\le0.01$ & $+2.58$ $[1.79,3.58]$ nats/step \\
 & network factor, chaotic $-$ $|\alpha|\le0.01$ & $-0.07$ $[-0.21,0.11]$ nats/step \\
 & entropy, chaotic $-$ periodic & $+1.62$ $[1.26,1.91]$ nats \\
 & top-2 logit gap$/T$, chaotic $-$ periodic & $-1.37$ $[-1.98,-0.87]$ \\
logit bias only, $\alpha=0$ & chaotic share & 3.5\,\% (241 of 6837 cells) \\
residual steering, block 12 & chaotic share & 6.2\,\% (331 of 5353 cells) \\
crossover, 0.5B, 14 cells & readout effect, arm 10 $-$ arm 00 & $+0.81$ $[0.29,1.35]$ nats/step \\
 & network effect, arm 01 $-$ arm 00 & $-0.49$ $[-0.63,-0.35]$ nats/step \\
 & interaction & $+0.26$ $[0.10,0.45]$ nats/step \\
crossover, 7B, 10 cells & readout effect & $+0.23$ $[-0.11,0.63]$ nats/step \\
 & network effect & $-1.14$ $[-1.61,-0.64]$ nats/step \\
\bottomrule
\end{tabular}
\caption{Where steering's instability comes from (Qwen2.5-0.5B unless marked 7B, Qwen2.5-7B). Intervals: bootstrap over
temperature rows (gain split) or over held-out cells (crossover). Runs MECH and XO; \texttt{tools/steer\_gain.py}, \texttt{tools/crossover.py},
\texttt{tools/plot\_mechanism.py}.}
\label{tab:mechanism}
\end{table}

